% !TEX root = ../notes.tex

\documentclass[../notes.tex]{subfiles}

\begin{document}

\section{May 11}
Here we go.

\subsection{Simpson's Conjecture}
Simpson's conjecture is about rigid local systems. Here is the definition.
\begin{notation}
	Fix a smooth projective variety $X$ over $\CC$. Let $\mc M_{\mathrm{Betti}}$ be the coarse moduli space of local systems on $X(\CC)$.
\end{notation}
\begin{remark}
	Equivalently, $\mc M_{\mathrm{Betti}}$ is the coarse moduli space of representations of $\pi_1(\CC)$. 
\end{remark}
\begin{remark}
	To make $M_{\mathrm{Betti}}$ well-behaved, one should only consider its components of fixed dimension.
\end{remark}
\begin{definition}[rigid]
	Fix a smooth projective variety $X$ over $\CC$. A \textit{rigid} local system is an isolated point of $M_{\mathrm{Betti}}(X)$ which is also semisimple.
\end{definition}
\begin{example}
	If $X$ is a complete curve of genus bigger than $1$, then $\mc M_{\mathrm{Betti}}$ is usually connected for fixed rank, so there are no rigid local systems. One sees that it is connected by doing some tangent space calculations; alternatively, one can show that it is a torsor of some group scheme over $\op{Bun}_n^\circ(X)$.
\end{example}
Here is a more flexible notion of rigid.
\begin{notation}
	Fix a smooth projective variety $X$ over $\CC$. Further, choose a nonnegative integer $n$ and a line bundle $\mc L$ on $X$. Let $\mc M_{\mathrm{Betti},n}^{\mc L}$ be the coarse moduli space of local systems on $X(\CC)$ of rank $n$ and a choice of isomorphism from the determinant to $\mc L$.
\end{notation}
\begin{definition}[rigid]
	Fix a smooth projective variety $X$ over $\CC$. Further, choose a nonnegative integer $n$ and a line bundle $\mc L$ on $X$. Then a \textit{rigid} local system is an isolated point on $\mc M_{\mathrm{Betti},n}^{\mc L}$ which is also semisimple.
\end{definition}
Here is our conjecture.
\begin{conj}[Simpson] \label{conj:simpson}
	Fix a smooth projective variety $X$ over $\CC$, and choose a line bundle $\mc L$ such that $\mc L^{\otimes d}=\OO_X$ for some positive $d$. Then a rigid local system $\mc E$ of determinant $\mc L$ is of geometric origin: $\mc E$ is a direct summand of a Gauss--Manin connection.
\end{conj}
\begin{remark}[Etingof]
	To get an interesting conjecture in rank one, one should allow from some ``punctures'' in codimension one. This is a ramified variant of this conjecture. It is known for $\PP^1$, due to Katz.
\end{remark}
Here is an application of the conjecture, which is known.
\begin{theorem} \label{thm:rigid-to-nilp}
	Fix a smooth projective variety $X$ over a number field $K$, and fix a rigid local system $\mc E$ with flat connection $\nabla$. Then for almost all primes $\mf p$ of $\OO_K$, the base-change $(\mc E_{k(\mf p)},\nabla_{k(\mf p)})$ has nilpotent $p$-curvature.
\end{theorem}
\Cref{conj:simpson} also makes it look like \Cref{conj:p-curv} may be more accessible for rigid local systems. As in our discussion of Katz's theorem, one can imagine proving \Cref{conj:p-curv} in two steps: unitarity and integrality. We are able to show one of these steps.
\begin{theorem} \label{thm:rigid-to-unitary}
	Fix a smooth projective variety $X$ over a number field $K$, and fix a rigid local system $\mc E$ with flat connection $\nabla$. If $\op{curv}_p(\mc E_{k(\mf p)},\nabla_{k(\mf p)})=0$ for all $\mf p$, then the monodromy is unitary.
\end{theorem}
Of course, if we could also show integrality, then we could finish the proof of Grothendieck's conjecture in this case, but this is currently out of reach.

\subsection{Nonabelian Hodge Theory}
For a smooth projective variety $X$, there is a filtration
\[0\to\mathrm H^{01}(X)\to\mathrm H^1_{\mathrm{dR}}(X)\to\mathrm H^{10}(X)\to0.\]
In positive characteristic, there is also a conjugate filtration
\[0\to\mathrm H^1(\OO_{X^{(1)}})\to\mathrm H^1_{\mathrm{dR}}(X)\to\mathrm H^0(\Omega^1_{X^{(1)}})\to0\]
arising from truncating the isomorphism $\mathrm H^i(\Omega^\bullet_{X^{(1)}})=\mathrm H^i(X^{(1)})$.

One may be interested in building a variant of these constructions with non-constant coefficients. For example, there is a filtration
\[0\to\mathrm H^0\left(\Omega^1\right)\to\mathrm H^1_{\mathrm{dR}}(\OO^\times)\to\mathrm H^1_{\mathrm{Zar}}(\OO^\times)\to0.\]
The middle group can be thought of as parameterizing line bundles with flat connection $(\mc L,\nabla)$. If $X$ is in positive characteristic and can be lifted to $W_2(k)$, then there is a conjugate filtration
\[0\to\mathrm H^1(\OO^\times_{X^{(1)}})\to\mathrm H^1_{\mathrm{dR}}(X;\OO^\times)\to\mathrm H^0(\Omega^1_{X^{(1)}})\to0,\]
where now the second map is given by the $p$-curvature.

To encode these notions geometrically, we notice that filtrations can be handled categorically.
\begin{remark}
	The category of vector spaces with a filtration is equivalent to the category of $\mathbb G_m$-equivariant line bundles on $\AA^1$. Roughly speaking, the fiber over $1$ should be the original space, and the fiber over $0$ should be the associated graded.
\end{remark}
Here is the story for $\OO$.
\begin{example}
	Thus, the filtered vector space $\mathrm H^1_{\mathrm{dR}}(X/\CC)$ gives rise to a $\mathbb G_m$-equivariant line bundle on $\AA^1_\CC$.
\end{example}
\begin{remark}
	Over $\CC$, the vector space $\mathrm H^1_{\mathrm{dR}}(X/\CC)$ equipped with its de Rham filtration also has an opposite filtration induced by taking complex conjugation, so the $\mathbb G_m$-equivariant bundle extends to $\CP^1$.
\end{remark}
Here is the story for $\OO^\times$.
\begin{notation}
	Fix a smooth projective variety $X$ over $\CC$. We define $\mc M_{\mathrm{dR}}$ to be the coarse moduli space of vector bundles with flat connection.
\end{notation}
\begin{example}[twistor family]
	Observe that $\mathrm H^1_{\mathrm{dR}}(\OO^\times)$ is $\mc M_{\mathrm{dR},1}$, where the $1$ denotes rank. In this case, Simpson observes that there is a family $\widetilde{\mc M}$ of $\mathbb G_m$-equivariant line bundles (parameterized over $\AA^1$) such that $\widetilde{\mc M}_1=\mc M_{\mathrm{dR},1}$ and $\widetilde{\mc M}_0$ is the moduli space of Higgs bundles. Namely, $\widetilde{\mc M}_0$ classifies pairs $(\mc E,\varphi)$ where $\varphi\colon\mc E\to\mc E\otimes\Omega^1$ has that the map $(\varpi\land\varphi)\colon\mc E\to\mc E\otimes\Omega^2$ vanishes. It again turns out that $\widetilde{\mc M}$ extends to $\CP^1$.
\end{example}
\begin{remark}
	The spaces $\mc M_{\mathrm{dR}}$ and $\mc M_{\mathrm{Betti}}$ are analytically isomorphic over $\CC$, but they are not algebraically isomorphic. Indeed, already in rank one over a curve $X$, $\mc M_{\mathrm{dR}}$ admits a filtration
	\[0\to\Omega^1\to\mathrm H^1_{\mathrm{dR}}(\OO^\times)\to\op{Pic}^\circ(X)\to0,\]
	but $\mc M_{\mathrm{Betti}}=(\CC^\times)^{2g}$. Thus, we see that we are basically looking for an algebraic splitting of the filtration, which need not exist in higher genus. However, there is a splitting in positive characteristic.
\end{remark}
It turns out that even the twistor bundle on $\CP^1$ can trivialize.
\begin{theorem}
	Fix a smooth projective variety $X$ over $\CC$. Then given a semisimple local system $(\mc E,\nabla)$, one can find a harmonic metric and thus an associated Higgs bundle $(\mc E,\Phi)$. This trivializes the twistor bundle.
\end{theorem}
\begin{remark}
	The metric is constructed by solving some nonlinear partial differential equations.
\end{remark}
One can even describe the sections.
\begin{theorem}[Simpson] \label{thm:simpson}
	Fix a smooth projective variety $X$ over $\CC$, and let $\widetilde{\mc M}$ be the twistor family. Then $\CC^\times$-invariant sections of $\widetilde{\mc M}$ are in correspondence with complex variations of Hodge structures.
\end{theorem}
\begin{remark}
	There is also a story for higher rank $n$. Here, there is a natural map $\mc M_{\mathrm{dR}}\to\op{Bun}_n(X)$ given by $(\mc E,\nabla)\mapsto\mc E$, which recovers one filtration. (If $X$ is a curve, then because flatness is uninteresting, one finds that $\mc M_{\mathrm{dR}}$ is a torsor over $T^*\op{Bun}_n(X)$.) The conjugate filtration in positive characteristic now looks like
	\[\mc M_{\mathrm{dR}}\to\bigoplus_{r=1}^n\mathrm H^0\left(\op{Sym}^r\Omega^1_{X^{(1)}}\right),\]
	where we send $(\mc E,\nabla)$ to the characteristic polynomial of the $p$-curvature. As usual, one can use this to define a spectral cover $\widetilde X$ of $X$, and it turns out that the fibers are torsors over $\op{Pic}\widetilde X$.
\end{remark}

\subsection{Some Sketches}
Let's now sketch our theorems.
\begin{proof}[Sketch of \Cref{thm:rigid-to-nilp}]
	We proceed in steps.
	\begin{enumerate}
		\item Over $\CC$, one can use \Cref{thm:simpson} to put rigid bundles with flat connection in bijection with rigid Higgs fields, where rigid Higgs fields are isolated points on $\mc M_{\mathrm{Higgs}}$. It turns out that rigid Higgs fields are nilpotent: indeed, otherwise, we can take a dilation to produce a nontrivial determinant.
		\item By the usual spreading out arguments, for almost all primes $\mf p$, the base-change $(\mc E,\nabla)\mapsto(\mc E_{k(\mf p)},\nabla_{k(\mf p)})$ will continue to provide a bijection between rigid bundles with flat connection and rigid flat connections over $k(\mf p)$. Similar holds for Higgs fields.
		\item Now, suppose that $p$ is large. Fix a rigid Higgs field $(\mc E_{k(\mf p)},\Phi_{k(\mf p)})$ on $X_{k(\mf p)}$, which we know is nilpotent. Because it is nilpotent, $(\mc E_{k(\mf p)},\Phi_{k(\mf p)})$ is supported in a $p$-neighborhood of the zero section $X_{k(\mf p)}\subseteq T^*X_{k(\mf p)}$. Then the inverse Cartier transform produces a rigid bundle with flat connection on $X_{k(\mf p)}$. But by counting, we see that we must be getting all rigid bundles with flat connection in this manner. In other words, every rigid bundle reduces to one with nilpotent $p$-curvature.
		\qedhere
	\end{enumerate}
\end{proof}
\begin{proof}[Sketch of \Cref{thm:rigid-to-unitary}]
	Fix a rigid bundle $(\mc E,\nabla)$ with flat connection. By \Cref{thm:simpson}, it comes from a complex variation of Hodge structures. Following the proof of \Cref{thm:katz-p-curv}, one needs to check that $\op{gr}(\mc E,\nabla)$ is supported on the zero section. Namely, we should show that the Kodaira--Spencer class vanishes.
	
	To this end, we claim that there are infinitely many primes $p$ for which the inverse Cartier transform of $(\mc E_{k(\mf p)},\nabla_{k(\mf p)})$ is $\op{gr}\mc E_{k(\mf p)}$, which will complete the proof by the hypothesis on $(\mc E,\nabla)$. This claim is proven using a positive characteristic Riemann--Hilbert correspondence.
\end{proof}
\begin{remark}
	The ``infinitely many primes'' in the second paragraph only refers to a positive-density set because they must use the following theorem of Cassels: for any algebraic extension $F$ of $\QQ$ generated by some set $\{\alpha_1,\ldots,\alpha_n\}$, there is a set of primes $p$ with positive density for which there is a map $F\to\ZZ_p$ sending each $\alpha_i$ to a $p$-adic unit.
\end{remark}

\end{document}
